\section{Deletion Sort}
\label{sec:b-deletion-sort}
\source{Codeforces 2200B (Notion: Bitwise, 800)}{Easy}

\problemstatement{An array of $n\le10$ positive integers. While it is not
non-decreasing, you must delete one element of your choice. Minimise the
number of elements left when the game ends.}

\begin{patternbox}
The tags say bitmask brute force ($2^{10}$ subsets), but a one-line
invariant solves it: \emph{keep an inversion alive}.
\end{patternbox}

\subsection*{Approach and intuition}

If the array is already sorted, the game ends at once: answer $n$. Otherwise
there is an inversion $a_i>a_j$ with $i<j$. As long as more than two elements
remain, delete an element other than $a_i$ and $a_j$: the inversion survives,
so the array is still unsorted and the game continues. With two elements left
(exactly the inversion), delete one; a single element is sorted. Answer $1$.

\subsection*{Algorithm}
\begin{algorithmic}[1]
\State \Return ($a$ non-decreasing) ? $n$ : $1$
\end{algorithmic}

\dryrun{$[1,4,2,3]$: unsorted, answer $1$. $[6,7]$: sorted, answer $2$.
\checkmark}

\subsection*{C++17 implementation}
\cppfile{bitwise/06_deletion_sort.cpp}

\begin{complexitybox}
\textbf{Time:} $O(n)$. \textbf{Space:} $O(n)$.
\end{complexitybox}

\begin{pitfallbox}
\begin{itemize}
  \item Equal neighbours are allowed (``non-decreasing'').
  \item If you do write the bitmask brute force, remember the game stops as
        soon as the array is sorted --- you cannot keep deleting.
\end{itemize}
\end{pitfallbox}
